% TOP_PROBLEM: {"id": 4100001, "title": "Classification of Finite Metric Spaces and Combinatorics of Convex Polytopes", "category_id": 6, "category": {"id": 6, "name": "geometry", "display_name": "Geometry", "description": "Euclidean and non-Euclidean geometry, geometric structures.", "slug": "geometry", "order_index": 6, "created_at": "2026-07-31T15:26:25.671Z"}, "status": "open"}
% FIRST_POSTED: null
% ENTRY_KIND: statement_only
% Imported from: batch14.tex; UnsolvedMath ID 4100001
% Compile twice: lualatex 4100001.tex
\documentclass[11pt]{article}
\usepackage{fontspec}
\setmainfont{Latin Modern Roman}
\usepackage{amsmath,amssymb,mathtools,unicode-math}
\setmathfont{Latin Modern Math}
\usepackage{xurl,hyperref}
\hypersetup{hidelinks,pdftitle={UnsolvedMath Problem 4100001}}
\newcommand{\sref}[2]{\hyperlink{source:#1:#2}{[S#2]}}
\newcommand{\cataloguescope}[1]{\paragraph{Mathematical question.}#1}
\begin{document}
% UnsolvedMath ID 4100001; source catalogue code AMR-040-0001
\setcounter{section}{4100000}
\section{Classification of Finite Metric Spaces and Combinatorics of Convex Polytopes}\label{4100001}
{\small\sffamily UnsolvedMath ID 4100001 · Geometry}
\subsection{Definitions and mathematical statement}

\paragraph{Shared notation.}$\mathbb N=\{1,2,\ldots\}$, and $\mathbb Z,\mathbb Q,\mathbb R,\mathbb C$ denote the integers, rationals, reals and complex numbers. Any other conventions are specified locally.


Let \(X\) be a finite set of cardinality \(n\ge2\) and let \(\rho\) be a metric on \(X\): it is symmetric, \(\rho(x,y)>0\) for \(x\ne y\), \(\rho(x,x)=0\), and it satisfies the triangle inequality. Write \(\delta_x\in\mathbb R^X\) for the coordinate vector with value one at \(x\) and zero elsewhere, and put
\[V_0(X)=\left\{v\in\mathbb R^X:\sum_{x\in X}v_x=0\right\},\quad
R_{X,\rho}=\operatorname{conv}\left\{\frac{\delta_x-\delta_y}{\rho(x,y)}:x,y\in X,\ x\ne y\right\}.\]
Here \(\operatorname{conv}\) denotes convex hull. The fundamental polytope \(R_{X,\rho}\) has dimension \(n-1\). Its face poset consists of its faces ordered by inclusion; its combinatorial type is the isomorphism class of this poset. Its \(f\)-vector records the number \(f_k\) of \(k\)-dimensional faces for \(0\le k\le n-1\), with \(f_{n-1}=1\). Call two finite metric spaces similar in the combinatorial sense if their fundamental polytopes have isomorphic face posets; this use of similar does not mean a distance-preserving similarity.
For \(X=\{1,\ldots,n\}\), let \(\mathcal M_n\) be the cone of symmetric distance matrices with zero diagonal, positive off-diagonal entries and triangle inequalities. A combinatorial-type class is
\[\mathcal S_\tau=\{\rho\in\mathcal M_n:R_{X,\rho}\text{ has face-poset type }\tau\}.\]
An open or generic type has a class containing a relatively open subset of \(\mathcal M_n\); its open components are the connected components of the relative interior of that class. A metric has Euclidean type if it admits an isometric embedding into a Euclidean space, equivalently into a real Hilbert space.
\paragraph{Statement recorded in the supplied source.}
Investigate all five parts of the classification problem:
\begin{enumerate}
\item Describe the face poset and its \(f\)-vector directly from inequalities in the distances \(\rho(x,y)\), in particular by linear inequalities where possible.
\item For each cardinality \(n\), estimate the number of distinct combinatorial types and its growth as \(n\to\infty\); in particular estimate the number of open, generic types.
\item Give sufficient conditions on two finite metric spaces for combinatorial similarity.
\item Describe the types represented by metrics of Euclidean type. Does every combinatorial type have such a representative?
\item Describe the possible topology of the combinatorial-type stratification of \(\mathcal M_n\). Is this classification universal in the source's sense that arbitrary semialgebraic singularity and stratification types can occur in the topology of its equivalence classes, or are there restrictions on the topological types of the open components?
\end{enumerate} \sref{4100001}{2}
\subsection{Short English statement}
Classify finite metric spaces using the faces of their fundamental polytopes: derive the face structure from distances, count the types and generic types, find similarity criteria, determine which types occur for Euclidean metrics, and understand the topology and possible universality of the resulting classes.
\subsection{Sources}
\begin{description}
\item[\hypertarget{source:4100001:1}{S1}] ULAM.AI and the UnsolvedMath Contributors, UnsolvedMath: A Curated Collection of Open Mathematics Problems, record 4100001 (AMR-040-0001). CC BY 4.0. \url{https://huggingface.co/datasets/ulamai/UnsolvedMath}.
\item[\hypertarget{source:4100001:2}{S2}] A. M. Vershik, Classification of Finite Metric Spaces and Combinatorics of Convex Polytopes (2015), arXiv:1504.03601, Section 2, Definition 1 and Problem 1, pp. 1–3, including footnote 1. \url{https://arxiv.org/pdf/1504.03601}.
\end{description}
\label{end:4100001}
\end{document}
